El código inicial proporcionado por la IA no presentó errores sintácticos en Mermaid, pero incurrió en una discrepancia estructural respecto al modelo de datos del proyecto y al caso de uso \texttt{cu\_us10}:
\begin{itemize}[leftmargin=*,nosep]
\item \textbf{Incoherencia con el modelo ER en la dependencia de \texttt{GUIA\_TALLA}:} La IA generó el endpoint \texttt{GET /api/products/\{id\}/size-guide} y la consulta \texttt{SELECT ... FROM size\_guides WHERE product\_id = \{id\}}, vinculando la guía de tallas a cada producto individual. En el diseño arquitectónico del equipo (\texttt{er.mmd} y \texttt{clases.mmd}), la entidad \texttt{GUIA\_TALLA} depende de \texttt{CATEGORIA} (guía anatómica por tipo de prenda) y el endpoint oficial del glosario es \texttt{GET /api/categorias/idCategoria/medidas}.
\item \textbf{Invocación del método \texttt{GuiaTalla.obtenerTablaMedidas(idCategoria)}:} Se sustituyó la consulta SQL inventada por la consulta de repositorio \texttt{findGuiaTallaByCategoria(idCategoria)} y el método definido en la clase \texttt{GuiaTalla} (\texttt{clases.mmd}), que devuelve \texttt{GuiaTalla[] (pechoCm, cinturaCm, caderaCm)}.
\item \textbf{Orden de apertura del modal y bloques opcionales:} Se eliminó la rama \texttt{404 Not Found} innecesaria, se situó la apertura de la ventana modal tras recibir \texttt{200 OK (GuiaTalla[])}, se modelaron los bloques \texttt{opt seleccionarTalla(talla)} y \texttt{opt cerrarGuia()}, y se añadió \texttt{mirrorActors: false}.
\end{itemize}
